% Options for packages loaded elsewhere
% Options for packages loaded elsewhere
\PassOptionsToPackage{unicode}{hyperref}
\PassOptionsToPackage{hyphens}{url}
%
\documentclass[
  english,
  russian,
  12pt,
  a4paper,
  DIV=11,
  numbers=noendperiod]{scrreprt}
\usepackage{xcolor}
\usepackage{amsmath,amssymb}
\setcounter{secnumdepth}{5}
\usepackage{iftex}
\ifPDFTeX
  \usepackage[T1]{fontenc}
  \usepackage[utf8]{inputenc}
  \usepackage{textcomp} % provide euro and other symbols
\else % if luatex or xetex
  \usepackage{unicode-math} % this also loads fontspec
  \defaultfontfeatures{Scale=MatchLowercase}
  \defaultfontfeatures[\rmfamily]{Ligatures=TeX,Scale=1}
\fi
\usepackage{lmodern}
\ifPDFTeX\else
  % xetex/luatex font selection
\fi
% Use upquote if available, for straight quotes in verbatim environments
\IfFileExists{upquote.sty}{\usepackage{upquote}}{}
\IfFileExists{microtype.sty}{% use microtype if available
  \usepackage[]{microtype}
  \UseMicrotypeSet[protrusion]{basicmath} % disable protrusion for tt fonts
}{}
\usepackage{setspace}
% Make \paragraph and \subparagraph free-standing
\makeatletter
\ifx\paragraph\undefined\else
  \let\oldparagraph\paragraph
  \renewcommand{\paragraph}{
    \@ifstar
      \xxxParagraphStar
      \xxxParagraphNoStar
  }
  \newcommand{\xxxParagraphStar}[1]{\oldparagraph*{#1}\mbox{}}
  \newcommand{\xxxParagraphNoStar}[1]{\oldparagraph{#1}\mbox{}}
\fi
\ifx\subparagraph\undefined\else
  \let\oldsubparagraph\subparagraph
  \renewcommand{\subparagraph}{
    \@ifstar
      \xxxSubParagraphStar
      \xxxSubParagraphNoStar
  }
  \newcommand{\xxxSubParagraphStar}[1]{\oldsubparagraph*{#1}\mbox{}}
  \newcommand{\xxxSubParagraphNoStar}[1]{\oldsubparagraph{#1}\mbox{}}
\fi
\makeatother


\usepackage{longtable,booktabs,array}
\usepackage{calc} % for calculating minipage widths
% Correct order of tables after \paragraph or \subparagraph
\usepackage{etoolbox}
\makeatletter
\patchcmd\longtable{\par}{\if@noskipsec\mbox{}\fi\par}{}{}
\makeatother
% Allow footnotes in longtable head/foot
\IfFileExists{footnotehyper.sty}{\usepackage{footnotehyper}}{\usepackage{footnote}}
\makesavenoteenv{longtable}
\usepackage{graphicx}
\makeatletter
\newsavebox\pandoc@box
\newcommand*\pandocbounded[1]{% scales image to fit in text height/width
  \sbox\pandoc@box{#1}%
  \Gscale@div\@tempa{\textheight}{\dimexpr\ht\pandoc@box+\dp\pandoc@box\relax}%
  \Gscale@div\@tempb{\linewidth}{\wd\pandoc@box}%
  \ifdim\@tempb\p@<\@tempa\p@\let\@tempa\@tempb\fi% select the smaller of both
  \ifdim\@tempa\p@<\p@\scalebox{\@tempa}{\usebox\pandoc@box}%
  \else\usebox{\pandoc@box}%
  \fi%
}
% Set default figure placement to htbp
\def\fps@figure{htbp}
\makeatother



\ifLuaTeX
\usepackage[bidi=basic,provide=*]{babel}
\else
\usepackage[bidi=default,provide=*]{babel}
\fi
% get rid of language-specific shorthands (see #6817):
\let\LanguageShortHands\languageshorthands
\def\languageshorthands#1{}


\setlength{\emergencystretch}{3em} % prevent overfull lines

\providecommand{\tightlist}{%
  \setlength{\itemsep}{0pt}\setlength{\parskip}{0pt}}



 
\usepackage[backend=biber,langhook=extras,autolang=other*]{biblatex}
\addbibresource{bib/cite.bib}

\usepackage[]{csquotes}

\usepackage{indentfirst}
\usepackage{float}
\floatplacement{figure}{H}
\IfFileExists{plex-otf.sty}{
  %% Full TeXlive
  \IfPackageAtLeastTF{plex-otf.sty}{2024-12-06}{
    \usepackage[math,RM={Scale=0.94},SS={Scale=0.94},SScon={Scale=0.94},TT={Scale=MatchLowercase,FakeStretch=0.9},DefaultFeatures={Ligatures=Common}]{plex-otf}
  }{
    \usepackage[RM={Scale=0.94},SS={Scale=0.94},SScon={Scale=0.94},TT={Scale=MatchLowercase,FakeStretch=0.9},DefaultFeatures={Ligatures=Common}]{plex-otf}
  }
}{
  %% TinyTeX
  \usepackage{libertine}
}
\KOMAoption{captions}{tableheading}
\usepackage{indentfirst}
\makeatletter
\@ifpackageloaded{caption}{}{\usepackage{caption}}
\AtBeginDocument{%
\ifdefined\contentsname
  \renewcommand*\contentsname{Содержание}
\else
  \newcommand\contentsname{Содержание}
\fi
\ifdefined\listfigurename
  \renewcommand*\listfigurename{Список иллюстраций}
\else
  \newcommand\listfigurename{Список иллюстраций}
\fi
\ifdefined\listtablename
  \renewcommand*\listtablename{Список таблиц}
\else
  \newcommand\listtablename{Список таблиц}
\fi
\ifdefined\figurename
  \renewcommand*\figurename{Рисунок}
\else
  \newcommand\figurename{Рисунок}
\fi
\ifdefined\tablename
  \renewcommand*\tablename{Таблица}
\else
  \newcommand\tablename{Таблица}
\fi
}
\@ifpackageloaded{float}{}{\usepackage{float}}
\floatstyle{ruled}
\@ifundefined{c@chapter}{\newfloat{codelisting}{h}{lop}}{\newfloat{codelisting}{h}{lop}[chapter]}
\floatname{codelisting}{Список}
\newcommand*\listoflistings{\listof{codelisting}{Листинги}}
\makeatother
\makeatletter
\makeatother
\makeatletter
\@ifpackageloaded{caption}{}{\usepackage{caption}}
\@ifpackageloaded{subcaption}{}{\usepackage{subcaption}}
\makeatother
\usepackage{bookmark}
\IfFileExists{xurl.sty}{\usepackage{xurl}}{} % add URL line breaks if available
\urlstyle{same}
\hypersetup{
  pdftitle={Отчёт о лабораторной работе},
  pdfauthor={Бабаев Рахим},
  pdflang={ru-RU},
  hidelinks,
  pdfcreator={LaTeX via pandoc}}


\title{Отчёт о лабораторной работе}
\usepackage{etoolbox}
\makeatletter
\providecommand{\subtitle}[1]{% add subtitle to \maketitle
  \apptocmd{\@title}{\par {\large #1 \par}}{}{}
}
\makeatother
\subtitle{Лабораторная работа №2. Расчёт сети Fast Ethernet}
\author{Бабаев Рахим}
\date{}
\begin{document}
\maketitle

\renewcommand*\contentsname{Содержание}
{
\setcounter{tocdepth}{1}
\tableofcontents
}
\listoffigures
\listoftables

\setstretch{1.5}
\chapter{Цель
работы}\label{ux446ux435ux43bux44c-ux440ux430ux431ux43eux442ux44b}

Целью лабораторной работы является изучение принципов технологий
Ethernet и Fast Ethernet, а также практическое освоение методики оценки
работоспособности сети Fast Ethernet по первой и второй моделям расчёта.

\chapter{Задание}\label{ux437ux430ux434ux430ux43dux438ux435}

В соответствии с методическими указаниями требуется оценить
работоспособность 100‑мегабитной сети Fast Ethernet. Для этого
необходимо: 1. Рассмотреть заданную топологию сети. 2. Выполнить
проверку по первой модели Fast Ethernet. 3. Выполнить проверку по второй
модели Fast Ethernet. 4. Сделать вывод о работоспособности сети.

Так как номер варианта в задании отдельно не указан, ниже приведён
подробный расчёт для всех вариантов из таблицы задания. Для сдачи отчёта
достаточно оставить свой вариант и удалить остальные.

\chapter{Теоретическое
введение}\label{ux442ux435ux43eux440ux435ux442ux438ux447ux435ux441ux43aux43eux435-ux432ux432ux435ux434ux435ux43dux438ux435}

Ethernet представляет собой технологию локальных сетей, в которой
передача данных осуществляется кадрами. В варианте Fast Ethernet
скорость передачи составляет 100 Мбит/с, при этом сохраняются основные
принципы Ethernet, включая формат кадра и механизм доступа к среде
CSMA/CD.

При оценке работоспособности сети Fast Ethernet используются две модели.

Первая модель основана на правилах построения сети. Для конфигурации с
сегментами 100BASE-TX и двумя повторителями класса II предельно
допустимый диаметр домена коллизий составляет 205 м. Кроме того, длина
каждого сегмента на витой паре не должна превышать 100 м.

Вторая модель основана на расчёте времени двойного оборота сигнала. Для
витой пары категории 5 удельное время двойного оборота составляет 1,112
битовых интервала на метр. Для пары терминалов TX/FX учитывается
задержка 100 би, для каждого повторителя класса II с портами TX/FX ---
92 би. Сеть считается работоспособной, если суммарное время двойного
оборота на наихудшем пути не превышает 512 би.

\chapter{Выполнение лабораторной
работы}\label{ux432ux44bux43fux43eux43bux43dux435ux43dux438ux435-ux43bux430ux431ux43eux440ux430ux442ux43eux440ux43dux43eux439-ux440ux430ux431ux43eux442ux44b}

\section{Общая схема
расчёта}\label{ux43eux431ux449ux430ux44f-ux441ux445ux435ux43cux430-ux440ux430ux441ux447ux451ux442ux430}

\begin{enumerate}
\def\labelenumi{\arabic{enumi}.}
\item
  Сначала была рассмотрена топология сети. По схеме видно, что в сети
  присутствуют два повторителя класса II и шесть сегментов типа
  100BASE-TX. Наихудший путь для расчёта проходит между наиболее
  удалёнными узлами через оба повторителя.
\item
  Для проверки по первой модели был определён диаметр домена коллизий.
  Для этого брались самый длинный сегмент слева от первого повторителя,
  соединительный сегмент между повторителями и самый длинный сегмент
  справа от второго повторителя.
\item
  После этого была выполнена проверка по второй модели. В расчёт
  включались задержка пары терминалов, задержки двух повторителей класса
  II и задержка кабельных сегментов по формуле: {[} T = 100 + 92 + 92 +
  1\{,\}112 \cdot L, {]} где (L) --- суммарная длина наихудшего пути в
  метрах.
\item
  На последнем этапе по каждому варианту был сделан вывод. Если диаметр
  домена коллизий не превышал 205 м и время двойного оборота не
  превышало 512 би, сеть считалась работоспособной. В противном случае
  делался вывод о несоответствии требованиям Fast Ethernet.
\end{enumerate}

\section{Вариант 1}\label{ux432ux430ux440ux438ux430ux43dux442-1}

\begin{enumerate}
\def\labelenumi{\arabic{enumi}.}
\item
  Для варианта 1 были заданы следующие длины сегментов: 96 м, 92 м, 80
  м, 5 м, 97 м и 97 м. Наихудший путь проходит между узлом на сегменте 1
  и узлом на сегменте 5 или 6, поэтому длина пути составляет: {[} L = 96
  + 5 + 97 = 198 \text{ м}. {]}
\item
  Проверка по первой модели показала, что все сегменты имеют длину менее
  100 м, а диаметр домена коллизий равен 198 м. Это меньше предельно
  допустимых 205 м, следовательно, по первой модели сеть работоспособна.
\item
  Далее был выполнен расчёт по второй модели: {[} T = 100 + 92 + 92 +
  1\{,\}112 \cdot 198 = 504\{,\}176 \text{ би}. {]} Полученное значение
  меньше 512 би, поэтому по второй модели сеть также работоспособна.
\item
  Итог по варианту 1: сеть удовлетворяет требованиям первой и второй
  моделей и может считаться работоспособной.
\end{enumerate}

\section{Вариант 2}\label{ux432ux430ux440ux438ux430ux43dux442-2}

\begin{enumerate}
\def\labelenumi{\arabic{enumi}.}
\item
  Для варианта 2 длины сегментов составляют 95 м, 85 м, 85 м, 90 м, 90 м
  и 98 м. Наихудший путь проходит между самым удалённым левым и правым
  узлами: {[} L = 95 + 90 + 98 = 283 \text{ м}. {]}
\item
  При проверке по первой модели было установлено, что хотя длина каждого
  отдельного сегмента не превышает 100 м, диаметр домена коллизий равен
  283 м. Это больше допустимых 205 м, поэтому по первой модели сеть
  неработоспособна.
\item
  По второй модели время двойного оборота равно: {[} T = 100 + 92 + 92 +
  1\{,\}112 \cdot 283 = 598\{,\}696 \text{ би}. {]} Это значение
  превышает предел 512 би, следовательно, по второй модели сеть также
  неработоспособна.
\item
  Итог по варианту 2: сеть не удовлетворяет требованиям Fast Ethernet ни
  по первой, ни по второй модели.
\end{enumerate}

\section{Вариант 3}\label{ux432ux430ux440ux438ux430ux43dux442-3}

\begin{enumerate}
\def\labelenumi{\arabic{enumi}.}
\item
  Для варианта 3 заданы длины сегментов 60 м, 95 м, 10 м, 5 м, 90 м и
  100 м. Наихудший путь проходит от узла на сегменте 2 к узлу на
  сегменте 6: {[} L = 95 + 5 + 100 = 200 \text{ м}. {]}
\item
  Проверка по первой модели показывает, что все сегменты имеют
  допустимую длину, а диаметр домена коллизий равен 200 м. Так как 200 м
  меньше 205 м, сеть по первой модели работоспособна.
\item
  Проверка по второй модели даёт следующий результат: {[} T = 100 + 92 +
  92 + 1\{,\}112 \cdot 200 = 506\{,\}4 \text{ би}. {]} Значение не
  превышает 512 би, следовательно, по второй модели сеть также
  работоспособна.
\item
  Итог по варианту 3: сеть соответствует требованиям обеих моделей и
  является работоспособной.
\end{enumerate}

\section{Вариант 4}\label{ux432ux430ux440ux438ux430ux43dux442-4}

\begin{enumerate}
\def\labelenumi{\arabic{enumi}.}
\item
  Для варианта 4 длины сегментов равны 70 м, 65 м, 10 м, 4 м, 90 м и 80
  м. Наихудший путь проходит между сегментом 1 и сегментом 5: {[} L = 70
  + 4 + 90 = 164 \text{ м}. {]}
\item
  По первой модели все сегменты удовлетворяют ограничению по длине, а
  диаметр домена коллизий составляет 164 м. Это существенно меньше 205
  м, следовательно, сеть работоспособна.
\item
  По второй модели время двойного оборота составляет: {[} T = 100 + 92 +
  92 + 1\{,\}112 \cdot 164 = 466\{,\}368 \text{ би}. {]} Полученный
  результат меньше 512 би, поэтому сеть удовлетворяет и второй модели.
\item
  Итог по варианту 4: сеть является работоспособной как по первой, так и
  по второй модели.
\end{enumerate}

\section{Вариант 5}\label{ux432ux430ux440ux438ux430ux43dux442-5}

\begin{enumerate}
\def\labelenumi{\arabic{enumi}.}
\item
  Для варианта 5 заданы длины сегментов 60 м, 95 м, 10 м, 15 м, 90 м и
  100 м. Наихудший путь проходит между сегментом 2 и сегментом 6: {[} L
  = 95 + 15 + 100 = 210 \text{ м}. {]}
\item
  При проверке по первой модели было получено, что диаметр домена
  коллизий равен 210 м. Это превышает допустимое значение 205 м,
  следовательно, по первой модели сеть неработоспособна.
\item
  По второй модели выполняется расчёт: {[} T = 100 + 92 + 92 + 1\{,\}112
  \cdot 210 = 517\{,\}52 \text{ би}. {]} Это значение превышает 512 би,
  следовательно, и по второй модели сеть также не удовлетворяет
  требованиям.
\item
  Итог по варианту 5: сеть неработоспособна по обеим моделям и требует
  уменьшения длины наихудшего пути.
\end{enumerate}

\section{Вариант 6}\label{ux432ux430ux440ux438ux430ux43dux442-6}

\begin{enumerate}
\def\labelenumi{\arabic{enumi}.}
\item
  Для варианта 6 длины сегментов составляют 70 м, 98 м, 10 м, 9 м, 70 м
  и 100 м. Наихудший путь проходит между сегментом 2 и сегментом 6: {[}
  L = 98 + 9 + 100 = 207 \text{ м}. {]}
\item
  По первой модели диаметр домена коллизий равен 207 м. Это больше
  допустимого значения 205 м, следовательно, по первой модели сеть
  неработоспособна.
\item
  По второй модели время двойного оборота рассчитывается так: {[} T =
  100 + 92 + 92 + 1\{,\}112 \cdot 207 = 514\{,\}184 \text{ би}. {]}
  Полученное значение немного превышает допустимый предел 512 би,
  поэтому по второй модели сеть также неработоспособна.
\item
  Итог по варианту 6: сеть не удовлетворяет требованиям Fast Ethernet и
  требует сокращения длины наихудшего пути.
\end{enumerate}

\chapter{Выводы}\label{ux432ux44bux432ux43eux434ux44b}

В ходе лабораторной работы были изучены основные принципы технологий
Ethernet и Fast Ethernet, а также методы оценки работоспособности сети
на основе ограничений диаметра домена коллизий и времени двойного
оборота сигнала.

Выполненный расчёт показал, что не всякая конфигурация сети Fast
Ethernet с сегментами короче 100 м автоматически является допустимой.
Для корректной работы сети необходимо учитывать не только длины
отдельных сегментов, но и суммарную длину наихудшего пути между наиболее
удалёнными узлами.

По результатам расчётов установлено, что вариант 1, вариант 3 и вариант
4 удовлетворяют требованиям первой и второй моделей. Вариант 2, вариант
5 и вариант 6 не проходят проверку, так как превышают либо допустимый
диаметр домена коллизий, либо предельное время двойного оборота, либо
оба ограничения одновременно.

Таким образом, при проектировании сети Fast Ethernet необходимо
обязательно выполнять расчёт по обеим моделям, поскольку именно это
позволяет заранее определить работоспособность конфигурации и избежать
ошибок при построении сети.

\chapter{Список
литературы}\label{ux441ux43fux438ux441ux43eux43a-ux43bux438ux442ux435ux440ux430ux442ux443ux440ux44b}

\begin{enumerate}
\def\labelenumi{\arabic{enumi}.}
\tightlist
\item
  Королькова А. В., Кулябов Д. С. Сетевые технологии. Лабораторный
  практикум.
\item
  Стандарт IEEE 802.3u Fast Ethernet.
\end{enumerate}


\printbibliography



\end{document}
